\chapter{حساب المثلثات في المثلث القائم الزاوية}\label{ch:g9:trig}

في \omterm{def:g6:shapes:triangles}{المثلث القائم الزاوية}، متى عرفت زاوية حادّة واحدة، تعيّن
\emph{شكل} المثلث --- ولا يبقى متغيّرًا سوى حجمه. ومن ثمّ لا
تتعلّق نسب أضلاعه إلا بالزاوية: وهذه النسب هي \omterm{def:g9:trig:ratios}{جيب التمام}
\omterm{def:g9:trig:ratios}{والجيب} والظلّ. وهي، مع مبرهنة فيثاغورس، تتيح لك حساب كل
ضلع وكل زاوية في \omterm{def:g6:shapes:triangles}{مثلث قائم الزاوية} اِنطلاقًا من معلومات قليلة
جدًّا.

\section{مبرهنة فيثاغورس، من جديد}

\begin{theorem}[فيثاغورس]\label{thm:g9:trig:pythagoras}
في مثلث $ABC$ \omterm{def:g6:shapes:triangles}{قائم الزاوية} في $A$، يكون مربّع الوتر
\omterm{def:g1:addition:def}{مجموع} مربّعَي ضلعَي القائمة:
\[
BC^2 = AB^2 + AC^2 .
\]
وبالعكس، إذا كان $BC^2 = AB^2 + AC^2$ في مثلث، كان \omterm{def:g6:shapes:triangles}{المثلث
قائم الزاوية} في $A$.
\end{theorem}
\admitted

\begin{example}\label{ex:g9:trig:pythagoras}
\omterm{def:g6:shapes:triangles}{مثلث قائم الزاوية} ضلعا قائمته $AB = 5$ و $AC = 12$. إذن
\[
BC^2 = 5^2 + 12^2 = 25 + 144 = 169,
\qquad\text{إذن}\quad BC = \sqrt{169} = 13 .
\]
ولإيجاد \emph{ضلع قائمة}: إذا كان $BC = 10$ و $AB = 6$، فإنّ
$AC^2 = BC^2 - AB^2 = 100 - 36 = 64$، إذن $AC = 8$.
\end{example}

\section{جيب التمام والجيب والظلّ}

في \omterm{def:g6:shapes:triangles}{المثلث القائم الزاوية}، وبالنسبة إلى زاوية حادّة $\theta$، تحمل
الأضلاع الثلاثة أسماء: \emph{الوتر} (المقابل \omterm{def:g3:shapes:rightangle}{للزاوية القائمة}، وهو
أطول الأضلاع)، والضلع \emph{المقابل} للزاوية $\theta$، والضلع
\emph{المجاور} للزاوية $\theta$ (وهو ضلع القائمة الملامس
للزاوية $\theta$).

\begin{definition}[النسب المثلثية]\label{def:g9:trig:ratios}
من أجل زاوية حادّة $\theta$ في \omterm{def:g6:shapes:triangles}{مثلث قائم الزاوية}:
\[
\cos\theta = \frac{\text{المجاور}}{\text{الوتر}},
\qquad
\sin\theta = \frac{\text{المقابل}}{\text{الوتر}},
\qquad
\tan\theta = \frac{\text{المقابل}}{\text{المجاور}} .
\]
وهذه النسب\index{جيب التمام}\index{جيب}\index{ظل} لا تتعلّق إلا
بالزاوية $\theta$، لا \omterm{def:g6:measure:volume}{بحجم} المثلث (فكل المثلثات القائمة الزاوية
التي لها الزاوية الحادّة نفسها تحجيمات بعضها لبعض، حسب طاليس).
\end{definition}

\begin{omfigure}
\begin{tikzpicture}[scale=1.15]
  \coordinate (A) at (0,0);
  \coordinate (B) at (4.2,0);
  \coordinate (C) at (4.2,2.2);
  \draw[thick] (A) -- (B) -- (C) -- cycle;
  \draw[thin] (B) ++(-0.28,0) -- ++(0,0.28) -- ++(0.28,0);
  \draw[omThm, ->] (0.9,0) arc (0:27.7:0.9);
  \node[omThm, font=\small] at (1.25,0.21) {$\theta$};
  \node[below, font=\small, omDef] at (2.1,0) {المجاور};
  \node[right, font=\small, omMeth] at (4.2,1.1) {المقابل};
  \node[above left=-2pt, font=\small, omProp] at (2.2,1.25) {الوتر};
  \node[below left, font=\small] at (A) {$B$};
  \node[below right, font=\small] at (B) {$A$};
  \node[above right, font=\small] at (C) {$C$};
\end{tikzpicture}

{\small الأضلاع الثلاثة منظورًا إليها من الزاوية $\theta = \widehat B$:
الوتر $[BC]$، والضلع المجاور $[BA]$، والضلع المقابل $[AC]$.}
\end{omfigure}

\begin{proposition}[الخواص الأولى]\label{prop:g9:trig:properties}
من أجل كل زاوية حادّة $\theta$:
\begin{enumerate}
  \item $0 < \cos\theta < 1$ و $0 < \sin\theta < 1$ (فضلع القائمة
        أقصر من الوتر)؛
  \item $\cos^2\theta + \sin^2\theta = 1$؛
  \item $\tan\theta = \dfrac{\sin\theta}{\cos\theta}$.
\end{enumerate}
\end{proposition}

\begin{proof}
لنكتب $a$ للضلع المجاور، و $o$ للضلع المقابل، و $h$ للوتر في
\omterm{def:g6:shapes:triangles}{مثلث قائم الزاوية} له الزاوية $\theta$.

1. $0 < a < h$ و $0 < o < h$، فالنسبتان محصورتان تمامًا بين
$0$ و $1$.

2. حسب فيثاغورس، $a^2 + o^2 = h^2$. فاِقسم كل شيء على $h^2$:
\[
\cos^2\theta + \sin^2\theta
= \frac{a^2}{h^2} + \frac{o^2}{h^2}
= \frac{a^2 + o^2}{h^2} = 1 .
\]

3. $\dfrac{\sin\theta}{\cos\theta} = \dfrac{o/h}{a/h} = \dfrac oa
= \tan\theta$.
\end{proof}

\begin{method}[إيجاد ضلع]\label{met:g9:trig:side}
لحساب ضلع مجهول حين يُعرف ضلع وزاوية حادّة:
\begin{enumerate}
  \item علّم الزاوية المعلومة وسمِّ الأضلاع الثلاثة (الوتر،
        والمقابل، والمجاور) \emph{اِنطلاقًا من تلك الزاوية}؛
  \item واِختر النسبة (\omterm{def:g9:trig:ratios}{جيب التمام} أو \omterm{def:g9:trig:ratios}{الجيب} أو الظلّ) التي تضمّ
        الضلع المعلوم والضلع المطلوب؛
  \item واُكتب \omterm{def:g8:equations:def}{المعادلة} وحُلّها؛
  \item وتحقّق من المعقولية: يجب أن يخرج الوتر أطول الأضلاع.
\end{enumerate}
\end{method}

\begin{example}\label{ex:g9:trig:side}
في \omterm{def:g6:shapes:triangles}{مثلث قائم الزاوية}، يقيس الوتر $8$ وإحدى الزاويتين
$35^\circ$؛ اِحسب الضلع المقابل لها. والنسبة التي تضمّ الضلع
المقابل والوتر هي \omterm{def:g9:trig:ratios}{الجيب}:
\[
\sin 35^\circ = \frac{\text{المقابل}}{8}
\quad\Longrightarrow\quad
\text{المقابل} = 8 \sin 35^\circ \approx 8 \times 0.574 \approx 4.6 .
\]
وللضلع المجاور، اِستعمل \omterm{def:g9:trig:ratios}{جيب التمام}:
$8\cos 35^\circ \approx 8 \times 0.819 \approx 6.6$. والتحقّق
بفيثاغورس: $4.6^2 + 6.6^2 \approx 21 + 43 \approx 64 = 8^2$.
\end{example}

\begin{method}[إيجاد زاوية]\label{met:g9:trig:angle}
حين يُعرف ضلعان، اِحسب النسبة التي يشكّلانها، وتعرّف عليها
جيبَ تمام أو جيبًا أو ظلًّا للزاوية المجهولة، وطبّق على
النسبة الدالة العكسية في الآلة الحاسبة ($\cos^{-1}$
أو $\sin^{-1}$ أو $\tan^{-1}$).
\end{method}

\begin{example}\label{ex:g9:trig:angle}
سلّم طوله $5$ م مسنود إلى جدار، وقدمه على بعد $1.4$ م من الجدار.
فالزاوية $\theta$ بين السلّم والأرض تحقّق
\[
\cos\theta = \frac{1.4}{5} = 0.28,
\qquad\text{إذن}\quad
\theta = \cos^{-1}(0.28) \approx 73.7^\circ .
\]
\end{example}

\begin{omfigure}
\begin{tikzpicture}[scale=0.85]
  \draw[thick] (0,0) -- (1.4,0);
  \draw[thick] (1.4,0) -- (1.4,4.0);
  \draw[omDef, very thick] (0,0) -- (1.4,4.0);
  \draw[thin] (1.4,0) ++(-0.25,0) -- ++(0,0.25) -- ++(0.25,0);
  \draw[omThm, ->] (0.75,0) arc (0:70.7:0.75);
  \node[omThm, font=\small] at (1.15,0.55) {$\theta$};
  \node[below, font=\small] at (0.7,0) {$1.4$ م};
  \node[omDef, font=\small, left] at (0.62,2.1) {$5$ م};
  \draw[gray, thin] (-0.7,0) -- (2.3,0);
\end{tikzpicture}

{\small مسألة السلّم: الضلع المجاور ($1.4$ م) والوتر ($5$ م)
معلومان، \omterm{def:g9:trig:ratios}{فجيب التمام} يعطي الزاوية.}
\end{omfigure}

\begin{example}[الزوايا الخاصّة]\label{ex:g9:trig:special}
يعطي مثلثان قيمًا مضبوطة. فنصف مربّع ضلعه 1 (وهو \omterm{def:g6:shapes:triangles}{مثلث
متساوي الساقين} \omterm{def:g6:shapes:triangles}{قائم الزاوية}) وتره $\sqrt2$ وزاويتاه $45^\circ$:
\[
\cos 45^\circ = \sin 45^\circ = \frac{1}{\sqrt2} = \frac{\sqrt2}{2},
\qquad \tan 45^\circ = 1 .
\]
ونصف \omterm{def:g6:shapes:triangles}{مثلث متساوي الأضلاع} ضلعه 1 يعطي الزاويتين $30^\circ$
و $60^\circ$:
\[
\cos 60^\circ = \sin 30^\circ = \frac12,
\qquad
\cos 30^\circ = \sin 60^\circ = \frac{\sqrt3}{2}.
\]
\end{example}

\section{تمارين}

\begin{exercise}[$\star$]\label{exo:g9:trig:1}
\omterm{def:g6:shapes:triangles}{مثلث قائم الزاوية} ضلعا قائمته $9$ و $12$. اِحسب وتره. وآخر وتره
$17$ وأحد ضلعَي قائمته $8$: اِحسب الضلع الآخر.
\end{exercise}

\begin{exercise}[$\star$]\label{exo:g9:trig:2}
مثلث أضلاعه $7$ و $24$ و $25$. فهل هو \omterm{def:g6:shapes:triangles}{قائم الزاوية}؟ والسؤال
نفسه من أجل الأضلاع $5$ و $6$ و $8$.
\end{exercise}

\begin{exercise}[$\star$]\label{exo:g9:trig:3}
في مثلث $DEF$ \omterm{def:g6:shapes:triangles}{قائم الزاوية} في $D$، وبتسمية الزاوية عند $E$
بالرمز $\theta$: أيّ ضلع هو الوتر؟ وأيّ ضلع هو المقابل للزاوية
$\theta$؟ واُكتب $\cos\theta$ و $\sin\theta$ و $\tan\theta$ نسبًا
بين الأضلاع $DE$ و $DF$ و $EF$.
\end{exercise}

\begin{exercise}[$\star$]\label{exo:g9:trig:4}
في \omterm{def:g6:shapes:triangles}{مثلث قائم الزاوية}، يقيس الوتر $10$ وإحدى الزاويتين الحادّتين
$28^\circ$. اِحسب ضلعَي القائمة (وقرّب إلى العُشر؛
و $\cos 28^\circ \approx 0.883$، و $\sin 28^\circ \approx 0.469$).
\end{exercise}

\begin{exercise}[$\star$]\label{exo:g9:trig:5}
في \omterm{def:g6:shapes:triangles}{مثلث قائم الزاوية}، يقيس ضلع القائمة المجاور للزاوية $\theta$
العدد $6$ ويقيس المقابل $4.5$. اِحسب $\tan\theta$، ثم $\theta$
(وقرّب إلى الدرجة؛ و $\tan^{-1}(0.75) \approx 36.9^\circ$).
\end{exercise}

\begin{exercise}[$\star\star$]\label{exo:g9:trig:6}
طائرة مسيّرة تطير على اِرتفاع $120$ م. ويراها مراقب على الأرض
بزاوية اِرتفاع قدرها $32^\circ$. فعلى أيّ مسافة أفقية من المراقب
تكون الطائرة
($\tan 32^\circ \approx 0.625$)؟
\end{exercise}

\begin{exercise}[$\star\star$]\label{exo:g9:trig:7}
منحدر ولوج يجب أن يرتفع $1.2$ م بزاوية لا تتجاوز $5^\circ$ مع
الأفق. فما أدنى طول على \omterm{prop:g9:linfunc:affinegraph}{الميل} يجب أن يكون للمنحدر
($\sin 5^\circ \approx 0.0872$)؟ قرّب بالزيادة إلى أقرب عُشر
متر.
\end{exercise}

\begin{exercise}[$\star\star$]\label{exo:g9:trig:8}
باِستعمال \cref{prop:g9:trig:properties}: زاوية حادّة $\theta$
تحقّق $\cos\theta = 0.6$. اِحسب $\sin\theta$ من غير إيجاد
$\theta$، ثم $\tan\theta$.
\end{exercise}

\begin{exercise}[$\star\star\star$]\label{exo:g9:trig:9}
برّر القيم المضبوطة في \cref{ex:g9:trig:special}: في \omterm{def:g6:shapes:triangles}{المثلث
المتساوي الساقين} \omterm{def:g6:shapes:triangles}{القائم الزاوية} الذي ضلعا قائمته $1$، اِحسب
الوتر؛ وفي \omterm{def:g6:shapes:triangles}{المثلث المتساوي الأضلاع} ذي الضلع $1$، اِحسب
الاِرتفاع، واِستنتج \omterm{def:g9:trig:ratios}{جيب التمام} \omterm{def:g9:trig:ratios}{والجيب} للزاويتين $30^\circ$
و $60^\circ$.
\end{exercise}

\section{مسألة: قياس ما لا تستطيع بلوغه}

\begin{problem}[{مسألة نهاية الأسبوع --- طريقة المساحين ذات
المحطّتين للاِرتفاعات التي لا تُبلَغ، ودرجات اِنحدار الطرق،
والمسافة إلى الأفق، وسباق الظلّ نحو اللانهاية}]
\label{pb:g9:trig:1}
لا يبلغ أيّ شريط قياس قمّة جبل، ولا قمّة برج وراء نهر، ولا
الأفق في البحر --- لكنّ مقياس زوايا ونسب هذا الفصل الثلاث
تبلغها. وقلب هذه المسألة هو \emph{طريقة المحطّتين} الكلاسيكية
عند المساحين، وهي تحسب اِرتفاعًا من غير الاِقتراب من قاعدته
أبدًا؛ وحولها: لافتات الطرق، والسلالم، والدرج، وتكوّر الأرض،
ولمحة أولى عن دالة تنفجر نحو اللانهاية.

\medskip
\noindent\textbf{الجزء الأول --- محطّة واحدة.}
(القيم: $\tan 35^\circ \approx 0.700$،
و $\sin 35^\circ \approx 0.574$، و $\cos 35^\circ \approx 0.819$،
و $\tan 1^\circ \approx 0.0175$.)
\begin{enumerate}
  \item من نقطة على بعد $60$ م من قاعدة منارة، تُرى القمّة
        بزاوية اِرتفاع قدرها $35^\circ$ (مقيسة من الأفق، عند
        مستوى العين). فكم تعلو القمّة فوق مستوى العين؟
  \item وتعلن لافتة طريق اِنحدارًا قدره $12\,\%$: أي أنّ الطريق
        يرتفع $12$ م لكل $100$ م أفقية. فأيّ زاوية يصنعها
        الطريق مع الأفق؟ وأيّ نسبة اِنحدار مئوية لطريق بزاوية
        $45^\circ$؟
  \item والزاويتان الحادّتان في \omterm{def:g6:shapes:triangles}{المثلث القائم الزاوية}
        متتامّتان، والضلع المقابل لكل منهما هو المجاور للأخرى.
        اِستنتج متطابقتَي «التمام»:
        \[
        \sin(90^\circ - \theta) = \cos\theta,
        \qquad
        \cos(90^\circ - \theta) = \sin\theta .
        \]
  \item وسلّم طوله $8$ م مسنود بزاوية $60^\circ$ مع الأرض.
        باِستعمال القيم المضبوطة في \cref{ex:g9:trig:special}،
        أعطِ الاِرتفاع المبلوغ وبعد القدم عن الجدار --- بالضبط،
        ثم إلى السنتيمتر.
  \item وتحقّق من $\cos^2\theta + \sin^2\theta = 1$ عدديًّا من
        أجل $\theta = 35^\circ$ وبالضبط من أجل
        $\theta = 60^\circ$. ولماذا تكون هذه المتطابقة
        (\cref{prop:g9:trig:properties}) عادةً جيدة للتحقّق من
        عمل الآلة الحاسبة؟
\end{enumerate}

\medskip
\noindent\textbf{الجزء الثاني --- محطّتان.}
قمّة جبل مرئية، وقاعدته لا تُبلَغ. فمن نقطة $A$، اِرتفاع القمّة
$30^\circ$؛ وبالمشي $200$ م مباشرةً نحو الجبل، إلى $B$، يصير
الاِرتفاع $40^\circ$. ولنكتب $h$ لعلوّ القمّة فوق مستوى العين
و $x$ للمسافة الأفقية من $B$ إلى عمودي القمّة.
($\tan 30^\circ \approx 0.577$،
و $\tan 40^\circ \approx 0.839$.)
\begin{enumerate}[resume]
  \item عبّر عن $\tan 40^\circ$ و $\tan 30^\circ$ بدلالة $h$
        و $x$ و $x + 200$: أي \omterm{def:g8:equations:def}{معادلتان} بمجهولين.
  \item ومن كل \omterm{def:g8:equations:def}{معادلة} عبّر عن المسافة الأفقية بدلالة $h$،
        واِطرح، واِستنتج صيغة المحطّتين:
        \[
        h = \frac{200}
        {\ \dfrac{1}{\tan 30^\circ} - \dfrac{1}{\tan 40^\circ}\ } .
        \]
        واِحسب $h$ إلى أقرب عشرة أمتار.
  \item واِحسب $x$ أيضًا، وتحقّق من معقولية جوابيك أمام رؤية
        $30^\circ$ من النقطة $A$.
  \item وفي جملة واحدة: لماذا لا غنى عن طريقة المحطّتين ---
        وأيّ قياس واحد، مستحيل هنا، كانت ستطلبه طريقة المحطّة
        الواحدة في السؤال 1؟
  \item ويُرى برج بزاوية $30^\circ$، وبزاوية $60^\circ$ بعد
        المشي $200$ م نحوه. وباِستعمال
        $\tan 30^\circ = \frac{1}{\sqrt3}$
        و $\tan 60^\circ = \sqrt3$، اِحسب الاِرتفاع
        \emph{بالضبط} --- والجواب \omterm{def:g9:arith:divisor}{مضاعف} للعدد
        $\sqrt3$ --- ثم إلى المتر.
\end{enumerate}

\medskip
\noindent\textbf{الجزء الثالث --- إلى الأفق وصعودًا على الجدار.}
\begin{enumerate}[resume]
  \item كم يبعد الأفق؟ إذا وقفت وعيناك على اِرتفاع $h$ مترًا
        فوق أرض كروية نصف قطرها $R$، لامس خطّ نظرك الكرة:
        فيشكّل خطّ النظر، \omterm{def:g3:shapes:shapes}{ونصف القطر} إلى نقطة الملامسة،
        \omterm{def:g3:shapes:shapes}{ونصف القطر} إلى قدميك مثلثًا \omterm{def:g6:shapes:triangles}{قائم الزاوية}. بيّن
        بفيثاغورس (\cref{thm:g9:trig:pythagoras}) أنّ المسافة
        $d$ إلى الأفق تحقّق $d^2 = 2Rh + h^2 \approx 2Rh$،
        واِحسب $d$ من أجل عينين على $1.80$ م
        ($R = 6.37 \times 10^6$ م).
  \item والسؤال نفسه من قمّة برج اِرتفاعه $324$ م. (وقاعدة
        البحّارة التقريبية: الأفق بالكيلومترات نحو
        $3.6\sqrt{h}$ حيث $h$ بالأمتار --- فتحقّق من جوابيك
        أمامها.)
  \item وإبهامك، وعرضه نحو $2$ سم، ممدودًا على طول ذراعك، على
        بعد نحو $60$ سم من عينك: فأيّ زاوية يغطّي؟ ويغطّي البدر
        نحو $0.5^\circ$: فأيّ جزء من إبهامك يخفي القمر كلّه؟
  \item والدرج المريح قوائمه $17$ سم ونوائمه $29$ سم. فبأيّ
        زاوية يصعد؟ قارن بالزاوية $30^\circ$ في
        \cref{ex:g9:trig:special}، وبسلّم علّية حادّ بزاوية
        $60^\circ$: فأيّ قائمة يحتاج \emph{هو} على نائمة قدرها
        $29$ سم؟
  \item وسباق الظلّ: اِحسب (بالآلة الحاسبة)
        $\tan 10^\circ$ و $\tan 45^\circ$ و $\tan 80^\circ$
        و $\tan 89^\circ$. فماذا يحدث حين تقترب الزاوية من
        $90^\circ$، ولماذا (فكّر في الضلع المجاور)؟ فالجدار
        عمودي، ولا قيمة للنسبة، ويقفز التمثيل البياني نحو
        اللانهاية --- وهذا أول \emph{خطّ مقارب} لك، وهو مخلوق
        يُدرس عن كثب في مجلّد الثانوية.
\end{enumerate}
\end{problem}
